\chapter{Build System}
\label{ch:build-system}

OmniLaTeX uses a Python-based build orchestrator, \texttt{build.py}, to manage
compilation of all documents in the repository.  At 2060 lines, it replaces
traditional Makefiles with explicit support for multiple build modes,
incremental builds via source-hash caching, parallel compilation, and rich
terminal output.

\section{Overview}

The build system is designed around several principles:

\begin{itemize}
    \item \textbf{Explicit modes.}  Every build selects a mode
          (\textsf{dev}, \textsf{prod}, or \textsf{ultra}) that determines the
          number of \LaTeX{} passes, bibliography processing, and glossary
          generation.
    \item \textbf{Incremental by default.}  A source-hash cache in
          \texttt{build/build\_cache.json} skips unchanged documents, reducing
          rebuild times from minutes to seconds.
    \item \textbf{Parallel execution.}  Examples are built concurrently using a
          thread pool, with progress displayed via the \ctanpackage{rich}
          library or a TQDM fallback.
    \item \textbf{Success via PDF existence.}  The orchestrator considers a
          build successful when a PDF file is produced, not when
          \texttt{latexmk} returns exit code zero.  This accommodates
          \LaTeX{} warnings that do not indicate actual failures.
\end{itemize}

\section{Build Commands}

Table~\ref{tab:build-commands} lists all top-level commands accepted by
\texttt{build.py}.

\begin{table}[htbp]
    \centering
    \caption{Build orchestrator commands}
    \label{tab:build-commands}
    \begin{tabular}{@{} l p{8.5cm} @{}}
        \toprule
        \textbf{Command} & \textbf{Description} \\
        \midrule
        \texttt{build-all} &
            Alias for the default build.  Compiles the root document and all
            examples. \\[4pt]
        \texttt{build-examples} &
            Build all example projects in \texttt{examples/} concurrently. \\[4pt]
        \texttt{build-example} &
            Build one or more named examples.  Accepts positional arguments:
            \texttt{build.py build-example article thesis}. \\[4pt]
        \texttt{build-root} &
            Build only the root document (\texttt{main.tex} at the repository
            root). \\[4pt]
        \texttt{clean} &
            Full cleanup.  Removes the entire \texttt{build/} directory,
            \texttt{\_minted/} caches, and \texttt{svg-inkscape/} caches. \\[4pt]
        \texttt{clean-aux} &
            Remove auxiliary files (\texttt{.aux}, \texttt{.log}, \texttt{.toc},
            \texttt{.bbl}, etc.) while preserving output PDFs. \\[4pt]
        \texttt{list} &
            List all available examples with their status (built/not built). \\[4pt]
        \texttt{preflight} &
            Run preflight checks: verify \texttt{latexmk}, \texttt{biber},
            \texttt{bib2gls}, \texttt{pygmentize}, Inkscape, and Gnuplot are
            available on \texttt{PATH}. \\[4pt]
        \texttt{test} &
            Run the Python test suite (pytest-based).  Does not require a full
            \LaTeX{} installation for unit tests; integration tests require
            Docker. \\
        \bottomrule
    \end{tabular}
\end{table}

Usage examples:

\begin{minted}{bash}
python3 build.py build-all
python3 build.py build-example article presentation
python3 build.py clean
python3 build.py list
\end{minted}

\section{Global Flags}

The following flags apply to all build commands:

\begin{table}[htbp]
    \centering
    \caption{Global build flags}
    \label{tab:build-flags}
    \begin{tabular}{@{} l p{8.5cm} @{}}
        \toprule
        \textbf{Flag} & \textbf{Description} \\
        \midrule
        \texttt{--mode \{dev,prod,ultra\}} &
            Select the build mode.  Defaults to \textsf{dev}.
            See \Cref{sec:build-modes} for details. \\[4pt]
        \texttt{--force} &
            Ignore the source-hash cache and rebuild everything from scratch.
            Equivalent to passing \texttt{-g} to \texttt{latexmk}. \\[4pt]
        \texttt{--timings} &
            Print per-example wall-clock times after the build completes.
            Useful for identifying slow examples. \\[4pt]
        \texttt{--verbose} &
            Enable verbose output from \texttt{latexmk} and subprocesses.
            Can also be set via the \texttt{OMNILATEX\_VERBOSE=1} environment
            variable. \\[4pt]
        \texttt{--source-date-epoch} &
            Set the \texttt{SOURCE\_DATE\_EPOCH} environment variable for
            reproducible builds.  Accepts a Unix timestamp. \\
        \bottomrule
    \end{tabular}
\end{table}

\section{Build Modes}
\label{sec:build-modes}

Each build mode controls how many \LaTeX{} passes are executed and whether
auxiliary processors (Biber, \texttt{bib2gls}) are invoked.

\begin{table}[htbp]
    \centering
    \caption{Build mode comparison}
    \label{tab:build-modes}
    \begin{tabular}{@{} l c c c @{}}
        \toprule
        \textbf{Feature} & \textbf{dev} & \textbf{prod} & \textbf{ultra} \\
        \midrule
        \LaTeX{} passes        & 3        & Full     & 2        \\
        Biber (bibliography)   & Yes      & Yes      & No       \\
        bib2gls (glossary)     & Yes      & Yes      & No       \\
        Target use case        & Editing  & Release  & Draft    \\
        Typical wall time      & $\sim$5\,min  & $\sim$14\,min & $\sim$2\,min \\
        \bottomrule
    \end{tabular}
\end{table}

\subsection{Dev Mode}

The default mode.  Runs three \LaTeX{} passes with Biber and \texttt{bib2gls}
processing.  Three passes are sufficient to resolve most cross-references and
citations during active editing.  Use this mode for day-to-day development.

\begin{minted}{bash}
python3 build.py build-example article --mode dev
\end{minted}

\subsection{Prod Mode}

Runs \texttt{latexmk} in its default mode, which performs as many passes as
needed until all references, citations, and glossary entries are fully resolved.
This is the mode used for final publication builds.

\begin{minted}{bash}
python3 build.py build-all --mode prod
\end{minted}

\subsection{Ultra Mode}

A minimal two-pass build with no bibliography or glossary processing.  Designed
for quick draft previews where reference accuracy is not critical.

\begin{minted}{bash}
python3 build.py build-example thesis --mode ultra
\end{minted}

\section{Incremental Builds}
\label{sec:incremental-builds}

The build system avoids redundant compilation through source-hash caching.

\subsection{How It Works}

Before building an example, \texttt{build.py} computes a SHA-256 hash of all
\texttt{.tex}, \texttt{.bib}, \texttt{.sty}, and \texttt{.cls} files in the
example directory and its dependencies.  This hash is stored in
\texttt{build/build\_cache.json} alongside the example name.

On subsequent builds, the current hash is compared against the cached value:

\begin{itemize}
    \item \textbf{Cache hit} -- The hash matches.  The example is skipped
          entirely, and the existing PDF is reused.
    \item \textbf{Cache miss} -- The hash differs (or no entry exists).  The
          example is rebuilt, and the cache is updated with the new hash.
\end{itemize}

\subsection{Bypassing the Cache}

Pass \texttt{--force} to ignore the cache and rebuild all targets:

\begin{minted}{bash}
python3 build.py build-all --force
\end{minted}

This is equivalent to removing \texttt{build/build\_cache.json} manually.  The
\texttt{--force} flag also passes \texttt{-g} to \texttt{latexmk}, which forces
a full recompilation even if \texttt{latexmk}'s own dependency tracking
considers the target up to date.

\section{latexmk Configuration}

The \texttt{.latexmkrc} file at the repository root configures \texttt{latexmk}
behaviour for all OmniLaTeX documents:

\begin{itemize}
    \item \textbf{Engine.}  \hologo{LuaLaTeX} is the default and only
          supported engine.  OmniLaTeX relies on Lua\TeX{} features including
          \ctanpackage{fontspec}, \ctanpackage{unicode-math}, and
          \ctanpackage{luatexja} for CJK support.
    \item \textbf{Interaction mode.}  \texttt{-interaction=nonstopmode} is used
          for all builds.  Errors are reported in the log but do not pause
          compilation.
    \item \textbf{Output directory.}  PDFs and auxiliary files are written to
          \texttt{build/}, keeping the source tree clean.
    \item \textbf{Minted cache.}  The \texttt{\_minted/} directory is
          configured as a cache for \ctanpackage{minted} output to avoid
          redundant Pygments invocations.
\end{itemize}

The entrypoint script in the Docker image also reads \texttt{.latexmkrc} from
the workspace, so the same configuration applies in containerised builds.

\section{Parallel Builds}
\label{sec:parallel-builds}

Example builds execute concurrently using Python's
\texttt{ThreadPoolExecutor}.  The default concurrency depends on the
environment:

\begin{itemize}
    \item \textbf{CI environments} (GitHub Actions, GitLab CI): 4 concurrent
          workers.
    \item \textbf{Local builds}: \texttt{os.cpu\_count()} or 4, whichever is
          lower.
\end{itemize}

The job count can be overridden with the \texttt{--jobs} flag:

\begin{minted}{bash}
python3 build.py build-examples --jobs 8
\end{minted}

Progress is displayed using the \texttt{rich} library when available, showing a
live table with per-example status (building/done/failed).  If \texttt{rich}
is not installed, a simple TQDM progress bar is used as a fallback.

Thread safety is ensured by creating the \texttt{build/examples/} directory
early and isolating per-example output into separate subdirectories.

\section{Performance}
\label{sec:build-performance}

\subsection{Full Build Timing}

A complete \texttt{build-all} in \textsf{prod} mode compiles the root document
and all 43 examples in approximately 14 minutes on a modern machine.  The same
build in \textsf{ultra} mode completes in roughly 2 minutes.

\subsection{Per-Example Timings}

Use the \texttt{--timings} flag to see wall-clock time for each example:

\begin{minted}{bash}
python3 build.py build-all --mode prod --timings
\end{minted}

This prints a summary table after all builds complete, sorted by duration.
Examples using \ctanpackage{tikz} or \ctanpackage{pgfplots} computations are
typically the slowest, as the \LaTeX{} engine performs numerical calculations
in-document rather than delegating to an external tool.

\subsection{Performance Tips}

\begin{itemize}
    \item Use \textsf{dev} mode during editing -- three passes are sufficient
          for most cross-reference resolution.
    \item Use \textsf{ultra} mode for rapid iteration when bibliography and
          glossary accuracy is not needed.
    \item Leverage incremental builds: only modified examples are recompiled.
    \item For \ctanpackage{pgfplots} heavy documents, pre-compute data
          externally and use \texttt{\textbackslash pgfplotstableread} instead
          of in-document calculations.
    \item Increase \texttt{--jobs} on machines with many cores, but note that
          each \hologo{LuaLaTeX} instance can consume significant memory.
\end{itemize}
